% Includable research insert. No preamble and no changes to the imported v1.
% Requires amsmath/amssymb/amsthm and lemma/theorem/proof environments.
% Bibliography keys shadrin, grosswendt, lin, gomes, gzyl are present in v1.
% The only hierarchy ingredient used below is the established weighted
% factor-32 squared-distance hierarchy theorem (v1, Lemma lem:hilbert).
% All logarithms are natural. Constants have not been optimized.

\section{A uniform all-capacity bound for polynomial curvature}
\label{sec:poly-allcap-new}

Let $\varphi\in C^2([0,1])$ have nonnegative, nonzero polynomial curvature
$g=\varphi''$ of degree at most an integer $r\ge1$, and write
\[
 B_\varphi(x,y)=\varphi(x)-\varphi(y)-\varphi'(y)(x-y),
 \qquad h(x)=\int_0^x\sqrt{g(t)}\,\mathrm dt,
 \qquad K_r=2\log(64r^2)+\frac14.
\]
The degree parameter counts curvature degree, not generator degree.
Since a nonzero polynomial cannot vanish on any interval, $h$ is continuous
and strictly increasing, and $\varphi$ is strictly convex.

\begin{lemma}[Polynomial curvature-coordinate comparison]
\label{lem:poly-h-new}
For every $x,y\in[0,1]$ and both orders of the arguments,
\begin{equation}
 \frac{1}{K_r}|h(x)-h(y)|^2
 \le B_\varphi(x,y)
 \le 8r^2|h(x)-h(y)|^2.
 \label{eq:poly-h-new}
\end{equation}
The constants are uniform over the polynomial coefficients and over the
location and length of the interval between $x$ and $y$.
\end{lemma}

\begin{proof}
The assertion is immediate for $x=y$. Fix $u<v$, and put
$L=v-u$, $\mu=\int_u^v g(t)\,\mathrm dt$, and
$M=\max_{[u,v]}g$. Both $\mu$ and $M$ are positive.
The Markov mass estimates used here are
\begin{equation}
 M\le\frac{8r^2\mu}{L},\qquad
 \int_u^v|t-a|g(t)\,\mathrm dt\ge\frac{L\mu}{64r^2}
 \quad (a\in[u,v]).
 \label{eq:poly-markov-new}
\end{equation}
For completeness, Markov's classical polynomial inequality
\cite{shadrin}, rescaled to $[u,v]$, gives
$\|g'\|_\infty\le2r^2M/L$. From a maximizer there is an available
one-sided segment of length $L/(4r^2)$, on which $g\ge M/2$.
Integration gives the first estimate. For any $a\in[u,v]$, the mass in
$\{t\in[u,v]:|t-a|<L/(32r^2)\}$ is at most
$LM/(16r^2)\le\mu/2$. The complementary mass gives the second estimate.

The two divergence orientations are
\[
 B_\varphi(v,u)=\int_u^v(v-t)g(t)\,\mathrm dt,
 \qquad
 B_\varphi(u,v)=\int_u^v(t-u)g(t)\,\mathrm dt.
\]
Fix either of these values and denote it by $B$. Measure the coordinate
$s\in[0,L]$ from the endpoint where its linear kernel vanishes, and denote
the resulting curvature by $\widetilde g(s)$. Explicitly,
$\widetilde g(s)=g(v-s)$ for $B_\varphi(v,u)$, and
$\widetilde g(s)=g(u+s)$ for $B_\varphi(u,v)$. Then
\[
 B=\int_0^L s\widetilde g(s)\,\mathrm ds,
 \qquad A:=h(v)-h(u)=\int_0^L\sqrt{\widetilde g(s)}\,\mathrm ds.
\]
Taking $a$ to be the zero-kernel endpoint in
\eqref{eq:poly-markov-new} gives
\begin{equation}
 \frac{L\mu}{64r^2}\le B\le L\mu,
 \qquad
 M\le\frac{512r^4B}{L^2}.
 \label{eq:poly-BM-new}
\end{equation}

Set $\delta=L/(64r^2)\in(0,L)$ and split $A=A_0+A_1$ at
$s=\delta$. The short endpoint contribution satisfies
\[
 A_0^2=\left(\int_0^\delta\sqrt{\widetilde g(s)}\,\mathrm ds\right)^2
 \le\delta^2M\le\frac{B}{8}.
\]
Weighted Cauchy--Schwarz gives
\[
 A_1^2
 =\left(\int_\delta^L
       \sqrt{s\widetilde g(s)}\,s^{-1/2}\,\mathrm ds\right)^2
 \le\left(\int_\delta^L s\widetilde g(s)\,\mathrm ds\right)
       \left(\int_\delta^L\frac{\mathrm ds}{s}\right)
 \le B\log(64r^2).
\]
Therefore $A^2\le2A_0^2+2A_1^2\le K_rB$, which proves the
left inequality in \eqref{eq:poly-h-new}.

For the other direction, $\sqrt{g(t)}\ge g(t)/\sqrt M$ implies
\[
 A^2\ge\frac{\mu^2}{M}\ge\frac{L\mu}{8r^2}.
\]
Together with $B\le L\mu$, this proves $B\le8r^2A^2$.
Every estimate applies to either orientation, completing the proof.
\end{proof}

\begin{theorem}[Polynomial all-capacity hierarchy bound]
\label{thm:poly-allcap-new}
Let $x_1,\ldots,x_N\in[0,1]$ have arbitrary positive weights
$p_1,\ldots,p_N$, and let $\varphi$ satisfy the preceding assumptions.
For a partition $P$ of $\{1,\ldots,N\}$ define
\[
 p_C=\sum_{i\in C}p_i,\qquad
 m_C=\frac{1}{p_C}\sum_{i\in C}p_ix_i,\qquad
 D_\varphi(P)=\sum_{C\in P}\sum_{i\in C}
          p_i\bigl(\varphi(x_i)-\varphi(m_C)\bigr).
\]
Let $\operatorname{OPT}_{\varphi,k}$ be the minimum of this cost over
partitions with at most $k$ cells. There is one deterministic hierarchy
$(H_k)_{k=1}^N$, with $H_{k+1}$ refining $H_k$ and $|H_k|\le k$, such that
\begin{equation}
 D_\varphi(H_k)
 \le 256r^2\left(2\log(64r^2)+\frac14\right)
           \operatorname{OPT}_{\varphi,k}
 \quad\text{for every }k\in\{1,\ldots,N\}.
 \label{eq:poly-allcap-new}
\end{equation}
In particular the all-capacity factor is
$O(r^2\log(r+2))$, uniformly in source size, positive weights, and posterior
locations. The statement includes capacities with zero flat optimum.
For normalized weights and the usual channel interpretation of distortion,
the same bound holds for stochastic all-capacity hierarchy price.

If $r=0$, positive constant curvature gives an all-capacity factor at most
$32$.
\end{theorem}

\begin{proof}
Put $y_i=h(x_i)$. For each cell $C$, the Bregman centroid identity gives
\begin{equation}
 D_\varphi(C)=\min_{q\in[0,1]}\sum_{i\in C}p_iB_\varphi(x_i,q),
 \label{eq:poly-centroid-new}
\end{equation}
because the expression on the right before minimization equals
$D_\varphi(C)+p_CB_\varphi(m_C,q)$. Define the transformed squared-distance
cost
\[
 V_h(P)=\sum_{C\in P}\sum_{i\in C}p_i(y_i-\overline y_C)^2,
 \qquad \overline y_C=\frac{1}{p_C}\sum_{i\in C}p_iy_i.
\]
The image $h([0,1])$ is an interval. It contains every
$\overline y_C$, so
\[
 \min_{q\in[0,1]}\sum_{i\in C}p_i|h(x_i)-h(q)|^2
 =\sum_{i\in C}p_i(y_i-\overline y_C)^2.
\]
Sum \eqref{eq:poly-h-new} with weights $p_i$, minimize over $q$ in each
cell, and use \eqref{eq:poly-centroid-new}. This gives for every partition
\begin{equation}
 \frac{1}{K_r}V_h(P)\le D_\varphi(P)\le8r^2V_h(P).
 \label{eq:poly-partition-new}
\end{equation}
This step does not require, and does not assert,
$h(m_C)=\overline y_C$.

Write $\operatorname{OPT}_{h,k}=\min_{|P|\le k}V_h(P)$.
The established weighted squared-distance hierarchy theorem supplies one
hierarchy with
$V_h(H_k)\le32\operatorname{OPT}_{h,k}$ at every capacity.
This is the positive-real-weight Hilbert-space transfer stated earlier;
its unweighted squared-Euclidean ingredient is recorded in
\cite[Lemma 2.13 and Theorem 1.2]{grosswendt} within the classical
hierarchy framework \cite{lin}.
From \eqref{eq:poly-partition-new},
$\operatorname{OPT}_{h,k}\le K_r\operatorname{OPT}_{\varphi,k}$.
Consequently
\[
 D_\varphi(H_k)
 \le8r^2V_h(H_k)
 \le256r^2\operatorname{OPT}_{h,k}
 \le256r^2K_r\operatorname{OPT}_{\varphi,k},
\]
as required. These are cost inequalities, so a zero optimum also forces
zero hierarchical cost; no division by an optimum is used.

For completeness, the stochastic flat optimum with output budget $k$
equals the deterministic flat optimum. Indeed, realize a channel by a
random deterministic function table $U$ independent of the input. Given
$U$, there are at most $k$ outputs. Conditional Jensen gives
\[
 \mathbb E\varphi(\mathbb E[x_X\mid Z])
 \le\mathbb E\varphi(\mathbb E[x_X\mid Z,U]),
\]
so the channel distortion is at least the average deterministic distortion
conditioned on $U$, and hence at least $\operatorname{OPT}_{\varphi,k}$.
Conversely any deterministic flat partition is a channel. The constructed
deterministic hierarchy is itself a stochastic degradation chain at all
capacities, giving the stochastic upper bound with the same denominators.

Finally, if $g\equiv a>0$ then
$B_\varphi(x,y)=a(x-y)^2/2$, so every partition cost is exactly
$a/2$ times ordinary squared-distance cost. The factor-$32$ theorem applies
directly, with no division by $r$.
\end{proof}

\paragraph{Scope and provenance.}
The comparison is a uniform quantitative specialization of the classical
curvature coordinate $h=\int\sqrt{\varphi''}$, which has antecedents in
separable Bregman geometry \cite{gomes,gzyl}. Its proof combines Markov's
inequality, an endpoint cutoff, and weighted Cauchy--Schwarz with an
established hierarchy transfer. No historical-priority claim is made.
The result improves the previously stated $O(r^4)$ all-capacity upper bound
for arbitrary finite sources and positive weights. It does not improve the
separate $O(r^2)$ bound for two prescribed capacities, prove a matching
all-capacity lower bound, or determine the exact optimal all-capacity order.
The lower construction in Section~\ref{sec:ext-weights} shows that a
quadratic degree order up to logarithmic factors is unavoidable, but does
not show that the logarithm in this upper bound is necessary. It makes no higher-dimensional
curvature assertion and gives an existence guarantee rather than a new
efficient optimization algorithm.
